\section{Radial incidence and convexity in an elementary fan}
\label{sec:area_law_fan_radial_incidence_convexity}

Let $o\in\R^2$, $\ell\in\N$ and $z\in\Z^2$ be arbitrary, and let
$Q=o+2^\ell(z+[0,1]^2)$ be the corresponding closed dyadic square.
Write $c$ for its center. On each side choose either the whole side or
its two midpoint halves. Orient the perimeter counterclockwise, and
orient each chosen base accordingly. Let $\mathcal E$ index the resulting
ordered bases $[a_i,b_i]$, and put
\begin{align}
    P_i & =\operatorname{conv}\{c,a_i,b_i\}, &
    R_i & =[c,b_i].
\end{align}
These are the triangles and radial segments of the elementary fan in
\cite[Section~11]{OpenAI2026AreaLaw}. All endpoints satisfy
$\|b_i-c\|_\infty=2^\ell/2>0$, and the directed rays determined by
$b_i-c$, for $i\in\mathcal E$, are pairwise distinct.

% Source: 10-geometry.tex, prop:two-families, lines 299--323,
% and geometry:initial-stars, lines 333--370;
% revision adc7f1241b42e322a6451854ab7e4b4c146bf78a.
% Manuscript: preprints/A-two-dimensional-area-law-from-a-global-spectral-gap-September-24-2026/
% build/sections/10-geometry.tex.
% The radial-incidence, affine-basis and radial-deletion auxiliaries have
% declaration references.
% Completion markers await ordinary-module and blueprint verification.

\begin{theorem}[Noncentral radial incidence]
    \label{thm:area_law_fan_noncentral_radial_incidence}
    \lean{TNLean.PEPS.AreaLaw.Geometry.cellFanPolygon_mem_iff_of_mem_radial}
    \uses{def:area_law_cell_fan}
    For $i,t\in\mathcal E$ and $x\in R_t$ with $x\ne c$,
    \begin{align}
        x\in P_i & \quad\Longleftrightarrow\quad
        t=i\ \text{or}\ b_t=a_i.
        \label{eq:fan_radial_incidence}
    \end{align}
\end{theorem}
\begin{proof}
    \uses{thm:area_law_fan_triangle_contact,
        thm:area_law_cell_fan_end_same_ray}
    The segment $R_t$ is contained in $P_t$ by convexity. If $t=i$,
    this proves the reverse implication. If $b_t=a_i$, then
    $R_t=[c,a_i]\subseteq P_i$, which gives the other reverse case.

    Conversely, suppose $x\in P_i$ and $t\ne i$. The distinct points
    $c,x$ belong to $P_t\cap P_i$. The classification of contacts
    between distinct fan triangles therefore gives either $b_t=a_i$,
    as required, or
    \begin{align}
        b_i & =a_t, & P_t\cap P_i & =R_i.
    \end{align}
    In the second case $x\in R_t\cap R_i$ and $x\ne c$. Thus there
    are $u,v>0$ with
    \begin{align}
        x-c & =u(b_t-c)=v(b_i-c).
    \end{align}
    The two endpoint vectors consequently determine the same directed
    ray. Distinctness of these directed rays gives $t=i$, contradicting
    the assumption. This proves~(\ref{eq:fan_radial_incidence}).
\end{proof}

\begin{theorem}[Affine basis of an elementary fan triangle]
    \label{thm:area_law_fan_triangle_affine_basis}
    \lean{TNLean.PEPS.AreaLaw.Geometry.cellFanPolygon_exists_affineBasis}
    \uses{def:area_law_cell_fan}
    For every $i\in\mathcal E$, the ordered triple $(c,a_i,b_i)$
    is an affine basis of $\R^2$.
\end{theorem}
\begin{proof}
    \uses{thm:area_law_cell_fan_regularity}
    The triangle $P_i$ has nonempty interior. Its affine span is
    therefore the whole plane. The vector space parallel to that
    affine span has dimension $2$. Since there are three vertices,
    they are affinely independent and form an affine basis.
\end{proof}

\begin{theorem}[Convexity after removal of a radial]
    \label{thm:area_law_fan_radial_deletion_convex}
    \lean{TNLean.PEPS.AreaLaw.Geometry.cellFanPolygon_sdiff_radial_convex}
    \uses{def:area_law_cell_fan}
    For every $i,t\in\mathcal E$, the set $P_i\setminus R_t$ is convex.
\end{theorem}
\begin{proof}
    \uses{thm:area_law_fan_noncentral_radial_incidence,
        thm:area_law_fan_triangle_affine_basis}
    Every $R_t$ contains $c$. Applying~(\ref{eq:fan_radial_incidence})
    away from $c$ therefore gives the exact alternatives
    \begin{align}
        P_i\cap R_t & =
        \begin{cases}
            [c,b_i], & t=i,                         \\
            [c,a_i], & t\ne i\ \text{and}\ b_t=a_i, \\
            \{c\},   & \text{otherwise}.
        \end{cases}
        \label{eq:fan_radial_face_cases}
    \end{align}
    The three vertices of $P_i$ form an affine basis by the preceding
    theorem. Let
    $\lambda_0,\lambda_1,\lambda_2$ be the barycentric coordinates for
    $c,a_i,b_i$, respectively. On $P_i$ these coordinates are
    nonnegative, sum to $1$, and reconstruct each point as
    $x=\lambda_0(x)c+\lambda_1(x)a_i+\lambda_2(x)b_i$.
    It follows that
    \begin{align}
        P_i\setminus[c,b_i]
         & =P_i\cap\{x:\lambda_1(x)>0\},              \\
        P_i\setminus[c,a_i]
         & =P_i\cap\{x:\lambda_2(x)>0\},              \\
        P_i\setminus\{c\}
         & =P_i\cap\{x:\lambda_1(x)+\lambda_2(x)>0\}.
        \label{eq:fan_radial_convex_deletions}
    \end{align}
    Each set in~(\ref{eq:fan_radial_convex_deletions}) is convex,
    being the intersection of a triangle with a strict affine
    half-plane. Hence~(\ref{eq:fan_radial_face_cases}) shows that
    $P_i\setminus R_t$ is convex for every $t\in\mathcal E$.
\end{proof}

\begin{theorem}[Convexity after removal of any family of radials]
    \label{thm:area_law_fan_radial_union_deletion_convex}
    \lean{TNLean.PEPS.AreaLaw.Geometry.cellFanPolygon_sdiff_radials_convex}
    \uses{def:area_law_cell_fan}
    For every $i\in\mathcal E$ and every subset $A\subseteq\mathcal E$,
    the set
    \begin{align}
        P_i\setminus\bigcup_{t\in A}R_t
        \label{eq:fan_radial_union_deletion}
    \end{align}
    is convex. The set $A$ may be empty.
\end{theorem}
\begin{proof}
    \uses{thm:area_law_fan_radial_deletion_convex}
    By the preceding theorem, $P_i\setminus R_t$ is convex for every
    $t\in\mathcal E$. Moreover,
    \begin{align}
        P_i\setminus\bigcup_{t\in A}R_t
         & =P_i\cap\bigcap_{t\in A}(P_i\setminus R_t).
    \end{align}
    The intersection of convex sets is convex. When $A=\varnothing$,
    the bounded intersection is the whole plane, so the last expression
    is $P_i$; no exclusion of the center is made in this case.
\end{proof}
